problem:  
Let \( (M^4, g_0) \) be a smooth, closed, oriented 4‑manifold with positive isotropic curvature, and assume the existence of a family of Ricci flows with surgery \( g_k(t) \), \(t>0\), with surgery scales \( \delta_k \to 0 \), such that \( (M, g_k(t)) \) is non‑collapsing for all \(k,t\).  
Suppose that, in the sense of **Bamler–Kleiner singular Ricci flow theory**, the space–time sequence \( (M,g_k(t)) \) converges (after passing to a subsequence) to a singular Ricci flow \( (\mathcal{X}, \{g_\infty(t)\}_{t>0}) \), where the regular subset \(\mathcal{X}_{\mathrm{reg}}\) carries a smooth Ricci flow and the singular set \(\mathcal{S} = \mathcal{X} \setminus \mathcal{X}_{\mathrm{reg}}\) has finite \(\mathcal{H}^3\)-measure.

Let \(\nu_k(x,\tau) := (4\pi\tau)^{-2} e^{-\ell_k(x,\tau)}\, dV_{g_k(\tau)}\) be the reduced‑density measures associated with each smooth \(g_k(t)\), where \(\ell_k\) is the reduced distance from a fixed basepoint \((p,0)\).  
Assume (as conjectured possible in this setting) that after passing to a subsequence, \(\nu_k(\cdot,\tau)\) converges weakly to a limit measure \(\nu_\infty(\tau)\) supported on \(\mathcal{X}_{\mathrm{reg}}\).  

**Question.** —  
1. Does there exist a natural extension of the pair \( (g_\infty(t), \nu_\infty(t)) \) as a **weak conjugate‑heat flow** across \(\mathcal{S}\), such that for all smooth test functions \(\phi\ge0\),
   \[
   \frac{d}{d\tau} \int_{\mathcal{X}} \phi\, d\nu_\infty(\tau)
   = -\int_{\mathcal{X}_{\mathrm{reg}}} \langle \nabla\phi,\nabla f_\infty\rangle e^{-f_\infty}(4\pi\tau)^{-2} dV_{g_\infty(\tau)} - \mu_{\mathrm{sing}}(\phi;\tau),
   \]
   where the defect term \(\mu_{\mathrm{sing}}(\cdot;\tau)\) is a non‑negative measure supported on \(\mathcal{S}\)?  

2. Is the total entropy dissipation at time \(\tau\),
   \[
   \mathcal{D}(\tau) := \int_{\mathcal{X}_{\mathrm{reg}}} \big| \mathrm{Ric}_{g_\infty(\tau)} + \nabla^2 f_\infty - \tfrac{1}{2\tau} g_\infty(\tau)\big|^2 e^{-f_\infty}(4\pi\tau)^{-2} dV_{g_\infty(\tau)} + \mu_{\mathrm{sing}}(\mathcal{X};\tau),
   \]
   finite for all \(\tau>0\)?  
   Equivalently, can the entire entropy dissipation of vanishing‑surgery limits be decomposed into a **smooth analytic part** and a **singular defect measure**, analogous to energy‑defect decompositions in other geometric evolution equations?

Why it is a "good" problem:  
This problem asks whether **entropy dissipation in 4‑dimensional singular Ricci flow limits admits a measure‑theoretic decomposition**, paralleling how “defect measures’’ arise in weak limits of critical geometric PDEs (harmonic map flow, Navier–Stokes, etc.).  
Unlike the monotonicity questions of previous attempts, this version captures a genuine unknown: whether the fundamental analytic invariant governing Ricci flow (Perelman’s entropy) retains a conservative structure that can be quantified even in the presence of singularities of co‑dimension one or higher.

If such a decomposition exists, it would define a canonical **entropy‑defect measure** quantifying “how much curvature concentration escapes” as one passes to the weak flow limit. This would bridge Ricci flow, geometric measure theory, and weak PDE analysis by introducing a new analytic invariant distinguishing smooth from singular evolutions.

If it fails, the precise mechanism of failure would expose fundamental obstructions to constructing 4‑dimensional singular Ricci flows—revealing the analytic “non‑closure’’ of Perelman’s entropy in higher dimensions.

Because the statement is concise, conceptually anchored in the known 3D theory, and analytically challenging—it probes the intersection of flow PDEs, measure convergence, and 4‑manifold singular structure—this question exemplifies a rigorous, high‑impact, “narrow entrance, profound depth’’ research‑level problem.